\section{BioDecisionBench: Design and Construction}
\label{sec:benchmark}

\subsection{Unified interface and record format}
\label{sec:schema}

Every example is one JSON line carrying
\texttt{task\_id}, \texttt{primitive}, \texttt{modality}, \texttt{sequences}
(each \texttt{\{role, modality, sequence\}}), \texttt{question}, \texttt{split},
\texttt{group}, and \texttt{provenance}, plus interface fields determined by the decision primitive
$\prim\in\{\text{noul},\text{choice},\text{score},\text{multi\_noul}\}$
(\cref{tab:primitives}). Sequences keep role boundaries for double-sequence tasks
(\texttt{query\_1}/\texttt{query\_2}: DNA+protein, DNA+DNA, protein+protein) rather than
being concatenated. \texttt{provenance} records source repository, revision/commit and
license; download receipts include per-file SHA-256.

\begin{table}[t]
\centering
\caption{The four decision primitives, their extra fields, supervision and metrics.}
\label{tab:primitives}
\small
\footnotesize
\begin{tabular}{p{0.10\textwidth}p{0.40\textwidth}p{0.36\textwidth}}
\toprule
Primitive & Extra fields & Supervision / evaluation \\
\midrule
choice & candidates (mutually exclusive labels), answer
       & CE; accuracy, macro-F1, NLL, candidate-validity rate \\
noul & candidates = \{no, yes\}, answer
     & binary CE; AUROC, MCC, accuracy, calibration \\
score & score\_levels (5), anchors (5 values from 4 interior train quantile
        cuts), gold\_value, gold\_level
      & soft-target CE + expected-value MSE; Spearman (primary), train-std-normalized
        RMSE \\
multi\_noul & candidates (all labels), per-label 0/1 vector, answer = positive set
            & masked binary; micro/macro-F1, per-label calibration \\
\bottomrule
\end{tabular}
\end{table}

\subsection{Sources, scale, and modalities}
\label{sec:scale}

\Cref{tab:benchmark} lists the composition. ProteinGym contributes 217 mutation-fitness
assays using the official contiguous-5 per-mutant folds (fold 0 = test, fold 1 = dev,
folds 2--4 = train; multi-mutant rows excluded); GUE 28 DNA choice tasks; RNAcompete 14
RNA score tasks (14-RBP subset, probes subsampled to 25k with a fixed seed); Genomic
Benchmarks 8; gene\_lan\_transfer 8 double-sequence noul tasks with endpoint
union-find component grouping; dnagpt DNA pools 3 (+1 promoter); TAPE 2; DeepSTARR 2
(dual outputs split into two tasks); DeepLoc 2.0 as one multi-noul task; and 7 local
snapshots from earlier public releases \citep{wang2024llamagene,wang2026biopaws}. Interface
distribution: choice 37 / noul 17 / score 235 / multi\_noul 1; modality: DNA 42 / protein
226 / RNA 14 / double-sequence 8; rows: 3{,}119{,}067 train / 440{,}305 dev /
581{,}242 test.

\emph{Counting convention.} The 290 figure counts evaluation units, not question types:
ProteinGym's 217 assays share one question type. We report both numbers wherever the
benchmark scale is quoted and never claim ``290 independent questions''.

% Table 1: benchmark composition (generated by gen_tables.py)
\begin{table}[t]
\centering
\caption{BioDecisionBench composition. Rows are the 10 source families (canonical
task-id prefix map used by the training scheduler; \texttt{lg\_}/\texttt{protein\_homology} share the
\texttt{local\_snapshots} family). Task counts are per evaluation unit; ProteinGym's 217 assays
share one question type. Admission: 0 train/test sequence leaks after isolation;
1{,}984 conflicting-label groups quarantined (not majority-voted);
925{,}985 cross-task duplicate sequences kept isolated per task.}
\label{tab:benchmark}
\small
\begin{tabular}{lrrrrrrrrr}
\toprule
Family & Mod. & choice & noul & score & multi & Tasks & Train & Dev & Test
\\ \midrule
ProteinGym & P & -- & -- & 217 & -- & 217 & 415,201 & 140,300 & 140,810 \\
GUE & D & 28 & -- & -- & -- & 28 & 731,069 & 82,350 & 82,125 \\
RNAcompete & R & -- & -- & 14 & -- & 14 & 275,449 & 33,610 & 33,476 \\
GenomicBenchmarks & D & 1 & 7 & -- & -- & 8 & 516,071 & 57,251 & 162,541 \\
gene\_lan & D+P & -- & 8 & -- & -- & 8 & 104,130 & 12,897 & 12,973 \\
local\_snapshots & D+P & 5 & 2 & -- & -- & 7 & 69,469 & 6,894 & 7,048 \\
dnagpt\_pools & D & 3 & -- & -- & -- & 3 & 111,269 & 12,527 & 13,892 \\
TAPE & P & -- & -- & 2 & -- & 2 & 74,955 & 7,874 & 40,042 \\
DeepSTARR & D & -- & -- & 2 & -- & 2 & 804,576 & 81,140 & 82,372 \\
DeepLoc & P & -- & -- & -- & 1 & 1 & 16,878 & 5,462 & 5,963 \\
\midrule
\textbf{Total} & D42+P226+R14+8D & 37 & 17 & 235 & 1 & 290 & 3,119,067 & 440,305 & 581,242 \\
\bottomrule
\end{tabular}
\end{table}


\subsection{Leak-proof admission protocol}
\label{sec:admission}

\begin{enumerate}
\item \textbf{Native splits first.} Sources with official train/dev(/test) splits are used
verbatim (GUE, Genomic Benchmarks, TAPE, DeepSTARR, ProteinGym folds, local snapshots).
\item \textbf{Legacy blind tests isolated from upstream pools.} The dnagpt pools (splice,
TF, core) ship as unsplit ``train'' data \emph{containing the blind-test members of prior
studies} \citep{wang2026biopaws}. Handling: test = legacy blind test, kept verbatim;
train/dev = the pool with \emph{all} legacy-test members removed, bucketed by sequence
hash; rows whose pool integer label conflicts with the blind-test text label (8 in core)
are quarantined as a group. Label mappings are not guessed but validated on full
pool$\cap$test overlaps (TF 3437/3437; core 5910/5918; splice agrees 4545/4545 under the
legacy convention 0$\to$Non-Splice/1$\to$Acceptor/2$\to$Donor, while the literature's newer
builder uses a different convention --- we adopt the data-consistent one and flag the
acceptor/donor semantics for independent verification).
\item \textbf{Exact cross-split sequence isolation.} A sequence appearing in several splits
is assigned wholly to the highest-priority split (test $>$ dev $>$ train); lower-split
copies are dropped. Same-sequence conflicting labels quarantine the whole group (1{,}984
groups; virus\_covid contributes 1{,}754) rather than majority-voting. We do
\emph{not} merge reverse complements: RC equivalence flips labels for strand-sensitive
tasks (splice donor/acceptor, promoter direction); label-aware RC grouping is left as
future work.
\item \textbf{Double-sequence tasks grouped by endpoints.} gene\_lan pairs are split by
union-find connected components over endpoints, so no endpoint appears in both train and
test. A construction-diagnosable subset (the two ``random'' pairing tasks are separable by
a frame-0 codon-table translation) is declared in the task notes and the paper, and is not
claimed as an independent blind test.
\item \textbf{Score anchor discipline.} The 5-level anchors are derived from training data
only (4 interior quantile cuts $\to$ 5 representative values); test values never feed back.
Raw RMSE is never averaged across assays (native scales differ by orders of magnitude);
Spearman is the primary metric.
\item \textbf{Admission checker.} Per task we verify split overlap, label conflicts,
answer $\in$ candidates, and length distributions; globally we count cross-task duplicate
sequences (925{,}985 sequence keys occur in $>$1 task --- not within-task leakage under a
multi-task model, but they must be escalated to global isolation under leave-one-task-out).
Final admission result: \textbf{0 within-task leaks, 0 flagged tasks}.
\end{enumerate}

\subsection{Scope boundaries}
\label{sec:boundaries}

Under the current encoder path (single sequence, DNA/protein, 512-token budget), 241/290
tasks are consumable by the compared architectures; the remaining 49 (14 RNA, 10
double-sequence, 1 multi-noul, 24 over-budget) are in the benchmark, covered by the
corresponding modality/interface evaluation paths or reported with exclusion rates ---
never silently truncated. The benchmark's value is the protocol and the data; it is not
bound to any single evaluated architecture.
